\section{AHSME 1986}
        \begin{enumerate}
        	\item (\textit{AHSME 1986}) $[x-(y-z)] - [(x-y) - z] =$


 $\textbf{(A)}\ 2y \qquad \textbf{(B)}\ 2z \qquad \textbf{(C)}\ -2y \qquad \textbf{(D)}\ -2z \qquad \textbf{(E)}\ 0$


 
 \textbf{Answer: }B



\textbf{Solution}

The expression becomes $(x-y+z)-(x-y-z) = x-y+z-x+y+z = 2z$, which is $\boxed{B}$.

 


	\item (\textit{AHSME 1986}) If the line $L$ in the $xy$-plane has half the slope and twice the $y$-intercept of the line $y = \frac{2}{3} x + 4$, then an equation for $L$ is:


 $\textbf{(A)}\ y = \frac{1}{3} x + 8 \qquad \textbf{(B)}\ y = \frac{4}{3} x + 2 \qquad \textbf{(C)}\ y =\frac{1}{3}x+4\qquad\\  \textbf{(D)}\ y =\frac{4}{3}x+4\qquad \textbf{(E)}\ y =\frac{1}{3}x+2$


 
 \textbf{Answer: }A



\textbf{Solution}

The original slope and $y$-intercept are $\frac{2}{3}$ and $4$, so the new ones are $\frac{1}{3}$ and $8$ respectively. Thus, using the slope-intercept form ($y = mx+c$), the new equation is $y=\frac{1}{3}x + 8$, which is $\boxed{A}$.

 


	\item (\textit{AHSME 1986}) $\triangle ABC$ has a right angle at $C$ and $\angle A = 20^\circ$. If $BD$ ($D$ in $\overline{AC}$) is the bisector of $\angle ABC$, then $\angle BDC =$


 $\textbf{(A)}\ 40^\circ \qquad \textbf{(B)}\ 45^\circ \qquad \textbf{(C)}\ 50^\circ \qquad \textbf{(D)}\ 55^\circ\qquad \textbf{(E)}\ 60^\circ$


 
 \textbf{Answer: }D



\textbf{Solution}

Since $\angle C = 90^{\circ}$ and $\angle A = 20^{\circ}$, we have $\angle ABC = 70^{\circ}$. Thus $\angle DBC = 35^{\circ}$. It follows that $\angle BDC = 90^{\circ} - 35^{\circ} = 55^{\circ}$, which is $\boxed{D}$.

 


	\item (\textit{AHSME 1986}) Let S be the statement 
"If the sum of the digits of the whole number $n$ is divisible by $6$, then $n$ is divisible by $6$."


 A value of $n$ which shows $S$ to be false is


 $\textbf{(A)}\ 30 \qquad \textbf{(B)}\ 33 \qquad \textbf{(C)}\ 40 \qquad \textbf{(D)}\ 42 \qquad \textbf{(E)}\ \text{ none of these}$


 
 \textbf{Answer: }B



\textbf{Solution}

For a counterexample, we need a number whose digit sum is divisible by $6$, but which is not itself divisible by $6$. $33$ satisfies these conditions, as $3+3=6$ but $6$ does not divide $33$, so the answer is $\boxed{B}$.

 


	\item (\textit{AHSME 1986}) Simplify $\left(\sqrt[6]{27} - \sqrt{6 \frac{3}{4} }\right)^2$


 $\textbf{(A)}\ \frac{3}{4} \qquad \textbf{(B)}\ \frac{\sqrt 3}{2} \qquad \textbf{(C)}\ \frac{3\sqrt 3}{4}\qquad \textbf{(D)}\ \frac{3}{2}\qquad \textbf{(E)}\ \frac{3\sqrt 3}{2}$


 
 \textbf{Answer: }A



\textbf{Solution}

We have $\sqrt[6]{27} = (3^{3})^{\frac{1}{6}} = 3^{\frac{1}{2}} = \sqrt{3}$ and $\sqrt{6 \frac{3}{4}} = \sqrt{\frac{27}{4}} = \frac{3 \sqrt{3}}{2}$, so the answer is $(\sqrt{3} - \frac{3 \sqrt{3}}{2})^{2} = (-\frac{\sqrt{3}}{2})^{2} = \frac{3}{4}$, which is $\boxed{A}$.

 


	\item (\textit{AHSME 1986}) Using a table of a certain height, two identical blocks of wood are placed as shown in Figure 1.  Length $r$ is found to be $32$ inches.  After rearranging the blocks as in Figure 2, length $s$ is found to be $28$ inches.  How high is the table?


 
                    \begin{figure}[H]
                        centering
                        \includegraphics[width=0.5\textwidth]{Problem/AHSME_I_1986_6_Problem_Figure_1.png}
                    \end{figure}
                    


 $\textbf{(A) }28\text{ inches}\qquad\textbf{(B) }29\text{ inches}\qquad\textbf{(C) }30\text{ inches}\qquad\textbf{(D) }31\text{ inches}\qquad\textbf{(E) }32\text{ inches}$


 
 \textbf{Answer: }C



\textbf{Solution}

Let $h$, $l$, and $w$ represent the height of the table and the length and width of the wood blocks, respectively, in inches. From Figure 1, we have $l+h-w=32$, and from Figure 2, $w+h-l=28$. Adding the equations gives $2h=60 \implies h=30$, which is $\boxed{C}$.

 


	\item (\textit{AHSME 1986}) The sum of the greatest integer less than or equal to $x$ and the least integer greater than or equal to $x$ is $5$. The solution set for $x$ is


 $\textbf{(A)}\ \Big\{\frac{5}{2}\Big\}\qquad \textbf{(B)}\ \big\{x\ |\ 2 \le x \le 3\big\}\qquad \textbf{(C)}\ \big\{x\ |\ 2\le x < 3\big\}\qquad\\  \textbf{(D)}\ \Big\{x\ |\ 2 < x\le 3\Big\}\qquad \textbf{(E)}\ \Big\{x\ |\ 2 < x < 3\Big\}$


 
 \textbf{Answer: }E



\textbf{Solution}

If $x \leq 2$, then $\lfloor x \rfloor + \lceil x \rceil \leq 2+2 < 5$, so there are no solutions with $x \leq 2$. If $x \geq 3$, then $\lfloor x \rfloor + \lceil x \rceil \geq 3+3$, so there are also no solutions here. Finally, if $2<x<3$, then $\lfloor x \rfloor + \lceil x \rceil = 2 + 3 = 5$, so the solution set is $\boxed{E}$.

 


	\item (\textit{AHSME 1986}) The population of the United States in $1980$ was $226,504,825$. The area of the country is $3,615,122$ square miles. There are $(5280)^{2}$ 
square feet in one square mile. Which number below best approximates the average number of square feet per person?


 $\textbf{(A)}\ 5,000\qquad \textbf{(B)}\ 10,000\qquad \textbf{(C)}\ 50,000\qquad \textbf{(D)}\ 100,000\qquad \textbf{(E)}\ 500,000$


 
 \textbf{Answer: }E



\textbf{Solution}

With about $230$ million people and under $4$ million square miles, there are about $60$ people per square mile. Since a square mile is about $(5000 \ \text{ft})^{2} = 25$ million square feet, that gives approximately $\frac{25}{60}$ of a million square feet per person. $\frac{25}{60}$ is approximately half, so the answer is approximately half a million, which is $\boxed{E}$.

 


	\item (\textit{AHSME 1986}) The product $\left(1-\frac{1}{2^{2}}\right)\left(1-\frac{1}{3^{2}}\right)\ldots\left(1-\frac{1}{9^{2}}\right)\left(1-\frac{1}{10^{2}}\right)$ equals


 $\textbf{(A)}\ \frac{5}{12}\qquad \textbf{(B)}\ \frac{1}{2}\qquad \textbf{(C)}\ \frac{11}{20}\qquad \textbf{(D)}\ \frac{2}{3}\qquad \textbf{(E)}\ \frac{7}{10}$


 
 \textbf{Answer: }C



\textbf{Solution 1}

Factor each term in the product as a difference of two squares, and group together all the terms that contain a $-$ sign, and all those that contain a $+$ sign. This gives $[(1-\frac{1}{2})(1-\frac{1}{3})(1-\frac{1}{4})...(1-\frac{1}{10})] \cdot [(1+\frac{1}{2})(1+\frac{1}{3})(1+\frac{1}{4})...(1+\frac{1}{10})] = [\frac{1}{2} \frac{2}{3} \frac{3}{4} ... \frac{9}{10}][\frac{3}{2} \frac{4}{3} \frac{5}{4} ... \frac{11}{10}] = \frac{1}{10} \cdot \frac{11}{2} = \frac{11}{20}$, which is $\boxed{C}$.

 



\textbf{Solution 2 (Answer Choices)}

Notice that the numerator of the last factor in this product is $99$, which is divisible by $11$.  Also, the denominators of the factors in this product are $2^2$, $3^2$, ..., $10^2$, none of which are divisible by $11$.  Therefore, the numerator of the answer must be divisible by $11$, and the only such answer is $\boxed{(\mathbf{C})\ \frac{11}{20}}$.

 -j314andrews

 


	\item (\textit{AHSME 1986}) The $120$ permutations of $AHSME$ are arranged in dictionary order as if each were an ordinary five-letter word. 
The last letter of the $86$th word in this list is:


 $\textbf{(A)}\ \text{A} \qquad \textbf{(B)}\ \text{H} \qquad \textbf{(C)}\ \text{S} \qquad \textbf{(D)}\ \text{M}\qquad \textbf{(E)}\ \text{E}$


 
 \textbf{Answer: }E



\textbf{Solution 1}

We could list out all of the possible combinations in dictionary order.

 $73rd:$ MAEHS                   $74th:$ MAESH$75th:$ MAHES$76th:$ MAHSE$77th:$ MASEH$78th:$ MASHE$79th:$ MEAHS$80th:$ MEASH$81th:$ MEHAS$82th:$ MEHSA$83th:$ MESAH$84th:$ MESHA$85th:$ MHAES$86th:$ MHASE

 We find that the $86th$ combination ends with the letter E. So the answer is $\textbf{(E)}\ E$.

 



\textbf{Solution 2}

We can do this problem without having to list out every single combination.There are $5$ distinct letters, so therefore there are $5!=120$ ways to rearrange the letters.We can divide the $120$ different combinations into 5 groups. Words that start with $A$, words that start with $E$ and so on...Combinations $1$-$24$ start with $A$,combinations $25$-$48$ start with $E$,combinations $49$-$72$ start with $H$,combinations $73$-$96$ start with $M$,and combinations $97$-$120$ start with $S$.We are only concerned with combination $86$, so we focus on combinations $73$-$96$.We can divide the remaining 24 combinations into 4 groups of 6, based upon the second letter.Combinations $73$-$78$ begin with $MA$,combinations $79$-$84$ begin with $ME$,combinations $85$-$90$ begin with $MH$,and combinations $91$-$96$ begin with $MS$.Combination $86$ begins with $MH$. Now we can fill in the rest of the letters in alphabetical order and get $MHASE$ (as $85$ is $MHAES$). The last letter of the word is $E$, so the answer is $\textbf{(E)}\ E$ .

 


	\item (\textit{AHSME 1986}) In $\triangle ABC, AB = 13, BC = 14$ and $CA = 15$. Also, $M$ is the midpoint of side $AB$ and $H$ is the foot of the altitude from $A$ to $BC$. 
The length of $HM$ is


 
                    \begin{figure}[H]
                        centering
                        \includegraphics[width=0.5\textwidth]{Problem/AHSME_I_1986_11_Problem_Figure_1.png}
                    \end{figure}
                    


 $\textbf{(A)}\ 6\qquad \textbf{(B)}\ 6.5\qquad \textbf{(C)}\ 7\qquad \textbf{(D)}\ 7.5\qquad \textbf{(E)}\ 8$


 
 \textbf{Answer: }B



\textbf{Solution}

In a right triangle, the length of the hypotenuse is twice the length of the median which bisects it. If the hypotenuse is $13$, the median must be $6.5$.

 



\textbf{Solution 2 (Self Torture)}

Warning: this solution is very intensive in calculation. Please do NOT try this on the test!

 Let's start by finding $AH$. By Heron's Formula, $s=\frac{13+14+15}{2}=21, [ABC]=\sqrt{21*(21-13)(21-14)(21-15)}=84$. Using the area formula $A=0.5bh$, $AH=12$. Now using the Pythagorean Theorem, $BH=5, HC=9$.

 Now $AM=MB=6.5$. Using Stewart's Theorem on $\triangle{ABH}$, letting $HM=x$:

 $13*6.5*6.5+13x^2=12*6.5*12+5*6.5*5$ (remember that Stewart's Theorem is $man+dad=bmb+cnc$).

 Thus $x=6.5$ or $x=-6.5$ (reject this solution since $x$ is positive). Thus $HM=6.5$. Select $\boxed{B}$.

 P.S.: Although this is torturous, this is a good practice of Heron's formula and Stewart's theorem though.

 


	\item (\textit{AHSME 1986}) John scores $93$ on this year's AHSME. Had the old scoring system still been in effect, he would score only $84$ for the same answers. 
How many questions does he leave unanswered? (In the new scoring system that year, one received $5$ points for each correct answer,
$0$ points for each wrong answer, and $2$ points for each problem left unanswered. In the previous scoring system, 
one started with $30$ points, received $4$ more for each correct answer, 
lost $1$ point for each wrong answer, and neither gained nor lost points for unanswered questions.)


 $\textbf{(A)}\ 6\qquad \textbf{(B)}\ 9\qquad \textbf{(C)}\ 11\qquad \textbf{(D)}\ 14\qquad \textbf{(E)}\ \text{Not uniquely determined}$


 
 \textbf{Answer: }B



\textbf{Solution}

Let $c$, $w$, and $u$ be the number of correct, wrong, and unanswered questions respectively. From the old scoring system, we have $30+4c-w=84$, from the new scoring system we have $5c+2u=93$, and since there are $30$ problems in the AHSME, $c+w+u=30$. Solving the simultaneous equations yields $u=9$, which is $\boxed{B}$.

 


	\item (\textit{AHSME 1986}) A parabola $y = ax^{2} + bx + c$ has vertex $(4,2)$. If $(2,0)$ is on the parabola, then $abc$ equals


 $\textbf{(A)}\ -12\qquad \textbf{(B)}\ -6\qquad \textbf{(C)}\ 0\qquad \textbf{(D)}\ 6\qquad \textbf{(E)}\ 12$


 
 \textbf{Answer: }E



\textbf{Solution}

Consider the quadratic in completed square form: it must be $y=a(x-4)^{2}+2$. Now substitute $x=2$ and $y=0$ to give $a=-\frac{1}{2}$. Now expanding gives $y=-\frac{1}{2}x^{2}+4x-6$, so the product is $-\frac{1}{2} \cdot 4 \cdot -6  = 3 \cdot 4 = 12$, which is $\boxed{E}$.

 


	\item (\textit{AHSME 1986}) Suppose hops, skips and jumps are specific units of length. If $b$ hops equals $c$ skips, $d$ jumps equals $e$ hops, 
and $f$ jumps equals $g$ meters, then one meter equals how many skips?


 $\textbf{(A)}\ \frac{bdg}{cef}\qquad \textbf{(B)}\ \frac{cdf}{beg}\qquad \textbf{(C)}\ \frac{cdg}{bef}\qquad \textbf{(D)}\ \frac{cef}{bdg}\qquad \textbf{(E)}\ \frac{ceg}{bdf}$


 
 \textbf{Answer: }D



\textbf{Solution}

$1$ metre equals $\frac{f}{g}$ jumps, which is $\frac{f}{g} \frac{e}{d}$ hops, and then $\frac{f}{g} \frac{e}{d} \frac{c}{b}$ skips, which becomes $\frac{cef}{bdg}$, i.e. answer $\boxed{D}$.

 


	\item (\textit{AHSME 1986}) A student attempted to compute the average, $A$, of $x, y$ and $z$ by computing the average of $x$ and $y$, and then computing the average of the result and $z$. Whenever $x < y < z$, the student's final result is


 $\textbf{(A)}\ \text{correct}\quad \textbf{(B)}\ \text{always less than A}\quad \textbf{(C)}\ \text{always greater than A}\quad\\ \textbf{(D)}\ \text{sometimes less than A and sometimes equal to A}\quad\\ \textbf{(E)}\ \text{sometimes greater than A and sometimes equal to A}  \quad$


 
 \textbf{Answer: }C



\textbf{Solution}

The true average is $A=\frac{x+y+z}{3}$, and the student computed $B=\frac{\frac{x+y}{2}+z}{2}=\frac{x+y+2z}{4}$, so $B-A = \frac{2z-x-y}{12} = \frac{(z-x)+(z-y)}{12}$, which is always positive as $z>x$ and $z>y$. Thus $B$ is always greater than $A$, i.e. $\boxed{C}$.

 


	\item (\textit{AHSME 1986}) In $\triangle ABC, AB = 8, BC = 7, CA = 6$ and side $BC$ is extended, as shown in the figure, to a point $P$ so that $\triangle PAB$
is similar to $\triangle PCA$. The length of $PC$ is


 
                    \begin{figure}[H]
                        centering
                        \includegraphics[width=0.5\textwidth]{Problem/AHSME_I_1986_16_Problem_Figure_1.png}
                    \end{figure}
                    


 $\textbf{(A)}\ 7\qquad \textbf{(B)}\ 8\qquad \textbf{(C)}\ 9\qquad \textbf{(D)}\ 10\qquad \textbf{(E)}\ 11$


 
 \textbf{Answer: }C



\textbf{Solution}

Since we are given that $\triangle{PAB}\sim\triangle{PCA}$, we have $\frac{PC}{PA}=\frac{6}{8}=\frac{PA}{PC+7}$.

 
Solving for $PA$ in $\frac{PC}{PA}=\frac{6}{8}=\frac{3}{4}$ gives us $PA=\frac{4PC}{3}$.

 
We also have $\frac{PA}{PC+7}=\frac{3}{4}$. Substituting $PA$ in for our expression yields $\frac{\frac{4PC}{3}}{PC+7}=\frac{3}{4}$

 
Which we can further simplify to $\frac{16PC}{3}=3PC+21$

 
$\frac{7PC}{3}=21$

 
$PC=9\implies\boxed{\textbf{(C)}}$

 


	\item (\textit{AHSME 1986}) A drawer in a darkened room contains $100$ red socks, $80$ green socks, $60$ blue socks and $40$ black socks. 
A youngster selects socks one at a time from the drawer but is unable to see the color of the socks drawn. 
What is the smallest number of socks that must be selected to guarantee that the selection contains at least $10$ pairs? 
(A pair of socks is two socks of the same color. No sock may be counted in more than one pair.)


 $\textbf{(A)}\ 21\qquad \textbf{(B)}\ 23\qquad \textbf{(C)}\ 24\qquad \textbf{(D)}\ 30\qquad \textbf{(E)}\ 50$


 
 \textbf{Answer: }B



\textbf{Solution 1}


Suppose that you wish to draw one pair of socks from the drawer. Then you would pick $5$ socks (one of each kind, plus one). Notice that in the worst possible situation, you will continue to draw the same sock, until you get $10$ pairs. This is because drawing the same sock results in a pair every $2$ of that sock, whereas drawing another sock creates another pair. Thus the answer is $5+2\cdot(10-1) = \boxed{\textbf{(B) } 23}$.

 



\textbf{Solution 2}

We have to choose $10$ $pairs$ of socks, i.e $20$ socks

 We choose $20$ socks in the worst possible scenario there will be $8$ $pairs$ and $4$ socks that are not paired, so we choose $4$ more socks to this, but $24$ can be the sum of $4$ $odd$ $numbers$ (like $7,7,7,3$)

 This means there might be $11$ $pairs$ and a few unpaired. If we choose $23$, it cannot be formed by summing $4$ $odd$ $numbers$. Therefore, $23$ must have $10$ $pairs$ of socks in it. Note that $21$ cannot happen because in the worst case scenario, we have $8$ $pairs$, and if $4$ $unpaired$, then pulling out one more will leave us with $9$ $pairs$. However, the question demands $10 pairs$. $22$ cannot happen either because if we pull out another sock it can be the same colour as the one we pulled out before, meaning it cannot be paired, but if we pull out another one and in the worst case scenario if the colour is same it still forms a pair with the last sock that we pulled out, therefore, the answer is $\boxed{\textbf{(B) } 23}$.

 



\textbf{Solution 3}

Since there are $4$ colors of socks, any selected set of socks could contain at most $4$ unpaired socks.  So a set of socks which contains $9$ or fewer matched pairs could have at most $4 + 2 \cdot 9 = 22$ socks.  This could happen if $19$ red socks and $1$ sock of each other color are selected.  So the minimum number of socks needed to guarantee $10$ matched pairs is $22 + 1 = \boxed{\textbf{(B) } 23}$.

 -j314andrews

 



\textbf{Solution 4}

Let $w$, $x$, $y$, and $z$ be the number of red, green, blue, and black socks selected, respectively.  Then the number of matched pairs of socks is $\left\lfloor \frac{w}{2} \right\rfloor + \left\lfloor \frac{x}{2} \right\rfloor + \left\lfloor \frac{y}{2} \right\rfloor + \left\lfloor \frac{z}{2} \right\rfloor$.  Note that for any integer $n$, $\left\lfloor \frac{n}{2} \right\rfloor = \frac{n}{2}$ if $n$ is even and $\left\lfloor \frac{n}{2} \right\rfloor = \frac{n-1}{2}$ if $n$ is odd.  So $\left\lfloor \frac{w}{2} \right\rfloor + \left\lfloor \frac{x}{2} \right\rfloor + \left\lfloor \frac{y}{2} \right\rfloor + \left\lfloor \frac{z}{2} \right\rfloor \geq \frac{w-1}{2} + \frac{x-1}{2} + \frac{y-1}{2} + \frac{z-1}{2} = \frac{w+x+y+z}{2} - 2$.  Thus if $\left\lfloor \frac{w}{2} \right\rfloor + \left\lfloor \frac{x}{2} \right\rfloor + \left\lfloor \frac{y}{2} \right\rfloor + \left\lfloor \frac{z}{2} \right\rfloor \leq 9$, then $\frac{w+x+y+z}{2} - 2 \leq 9$ and $w+x+y+z \leq 22$.  That is, a set of socks containing $9$ or fewer matched pairs contains at most $22$ socks.  Several equality cases exist, such as $(w, x, y, z) = (19, 1, 1, 1)$.  So the minimum number of socks needed to guarantee $10$ matched pairs is $22 + 1 = \boxed{\textbf{(B) } 23}$.

 -j314andrews

 


	\item (\textit{AHSME 1986}) A plane intersects a right circular cylinder of radius $1$ forming an ellipse. 
If the major axis of the ellipse is $50\%$ longer than the minor axis, the length of the major axis is


 $\textbf{(A)}\ 1\qquad \textbf{(B)}\ \frac{3}{2}\qquad \textbf{(C)}\ 2\qquad \textbf{(D)}\ \frac{9}{4}\qquad \textbf{(E)}\ 3$


 
 \textbf{Answer: }E



\textbf{Solution 1}

The length of the minor axis is the distance from the opposite points of the cylinder, which is simply $2 \cdot 1 =2$, and according to the question, the length of the major axis is simply $2 \cdot 150\% = \boxed{(E) 3}$

 



\textbf{Solution 2}

We note that we can draw the minor axis to see that because the minor axis is the minimum distance between two opposite points on the ellipse, we can draw a line through two opposite points of the cylinder, and so the minor axis is $2(1) = 2$. Therefore, our answer is $2(1.5) = 3$, and so our answer is $\boxed{E}$.

 


	\item (\textit{AHSME 1986}) A park is in the shape of a regular hexagon $2$ km on a side. Starting at a corner, 
Alice walks along the perimeter of the park for a distance of $5$ km. 
How many kilometers is she from her starting point?


 $\textbf{(A)}\ \sqrt{13}\qquad \textbf{(B)}\ \sqrt{14}\qquad \textbf{(C)}\ \sqrt{15}\qquad \textbf{(D)}\ \sqrt{16}\qquad \textbf{(E)}\ \sqrt{17}$


 
 \textbf{Answer: }A



\textbf{Solution}

We imagine this problem on a coordinate plane and let Alice's starting position be the origin. We see that she will travel along two edges and then go halfway along a third. Therefore, her new $x$-coordinate will be $1 + 2 + \frac{1}{2} = \frac{7}{2}$ because she travels along a distance of $2 \cdot \frac{1}{2} = 1$ km because of the side relationships of an equilateral triangle, then $2$ km because the line is parallel to the $x$-axis, and the remaining  distance is $\frac{1}{2}$ km because she went halfway along and because of the logic for the first part of her route. For her $y$-coordinate, we can use similar logic to find that the coordinate is $\sqrt{3} + 0 - \frac{\sqrt{3}}{2} =  \frac{\sqrt{3}}{2}$. Therefore, her distance is 
 \[\sqrt{\left(\frac{7}{2}\right)^2 + \left(\frac{\sqrt{3}}{2}\right)^2} = \sqrt{\frac{49}{4} + \frac{3}{4}} = \sqrt{\frac{52}{4}} = \sqrt{13},\] giving an answer of $\boxed{A}$.

 


	\item (\textit{AHSME 1986}) Suppose $x$ and $y$ are inversely proportional and positive. If $x$ increases by $p\%$, then $y$ decreases by 


 $\textbf{(A)}\ p\%\qquad \textbf{(B)}\ \frac{p}{1+p}\%\qquad \textbf{(C)}\ \frac{100}{p}\%\qquad \textbf{(D)}\ \frac{p}{100+p}\%\qquad \textbf{(E)}\ \frac{100p}{100+p}\%$


 
 \textbf{Answer: }E



\textbf{Solution}

We see that $x$ is multiplied by $\frac{100+p}{100}$ when it is increased by $p$\%. Therefore, $y$ is multiplied by $\frac{100}{100+p}$, and so it is decreased by $\frac{p}{100+p}$ times itself. Therefore, it is decreased by $\frac{100p}{100+p}$\%, and so the answer is $\boxed{E}$.

 


	\item (\textit{AHSME 1986}) In the configuration below, $\theta$ is measured in radians, $C$ is the center of the circle, 
$BCD$ and $ACE$ are line segments and $AB$ is tangent to the circle at $A$.


 
                    \begin{figure}[H]
                        centering
                        \includegraphics[width=0.5\textwidth]{Problem/AHSME_I_1986_21_Problem_Figure_1.png}
                    \end{figure}
                    


 A necessary and sufficient condition for the equality of the two shaded areas, given $0 < \theta < \frac{\pi}{2}$, is


 $\textbf{(A)}\ \tan \theta = \theta\qquad \textbf{(B)}\ \tan \theta = 2\theta\qquad \textbf{(C)}\ \tan\theta = 4\theta\qquad \textbf{(D)}\ \tan 2\theta =\theta\qquad\\  \textbf{(E)}\ \tan\frac{\theta}{2}=\theta$


 
 \textbf{Answer: }B



\textbf{Solution}

Well, the shaded sector's area is basically $\text{(ratio of } \theta \text{ to total angle of circle)} \times \text{(total area)} = \frac{\theta}{2\pi} \cdot (\pi r^2) = \frac{\theta}{2} \cdot (AC)^2$.

 In addition, if you let $\angle{ACB} = \theta$, then 
 \[\tan \theta = \frac{AB}{AC}\]
 \[AB = AC\tan \theta = r\tan \theta\]
 \[[ABC] = \frac{AB \cdot AC}{2} = \frac{r^2\tan \theta}{2}\]Then the area of that shaded thing on the left becomes 
 \[\frac{r^2\tan \theta}{2} - \frac{\theta \cdot r^2}{2}\]We want this to be equal to the sector area so 
 \[\frac{r^2\tan \theta}{2} - \frac{\theta \cdot r^2}{2} = \frac{\theta \cdot r^2}{2}\]
 \[\frac{r^2\tan \theta}{2} = \theta \cdot r^2\]
 \[\tan \theta = 2\theta\]

 


	\item (\textit{AHSME 1986}) Six distinct integers are picked at random from $\{1,2,3,\ldots,10\}$. What is the probability that, among those selected, the second smallest is $3$?


 $\textbf{(A)}\ \frac{1}{60}\qquad \textbf{(B)}\ \frac{1}{6}\qquad \textbf{(C)}\ \frac{1}{3}\qquad \textbf{(D)}\ \frac{1}{2}\qquad \textbf{(E)}\ \text{none of these}$


 
 \textbf{Answer: }C



\textbf{Solution}

The total number of ways to choose 6 numbers is ${10\choose 6} = 210$.

 Assume $3$ is the second-lowest number. There are $5$ numbers left to choose, $4$ of which must  be greater than $3$, and $1$ of which must be less than $3$. This is equivalent to choosing $4$ numbers from the $7$ numbers larger than $3$, and $1$ number from the $2$ numbers less than $3$.  
 \[{7\choose 4} {2\choose 1}= 35\times2.\]

 Thus, $\frac{35\times2}{210} = \frac{1}{3}$, which is $\fbox{C}$.

 


	\item (\textit{AHSME 1986}) Let $N = 69^{5} + 5 \cdot 69^{4} + 10 \cdot 69^{3} + 10 \cdot 69^{2} + 5 \cdot 69 + 1$. How many positive integers are factors of $N$?


 $\textbf{(A)}\ 3\qquad \textbf{(B)}\ 5\qquad \textbf{(C)}\ 69\qquad \textbf{(D)}\ 125\qquad \textbf{(E)}\ 216$


 
 \textbf{Answer: }E



\textbf{Solution}

Let $69=a$. Therefore, the equation becomes $a^5+5a^4+10a^3+10a^2+5a+1$. From Pascal's Triangle, we know this equation is equal to $(a+1)^5$. Simplifying, we have the desired sum is equal to $70^5$ which can be prime factorized as $2^5\cdot5^5\cdot7^5$. Finally, we can count the number of factors of this number.

 $6\cdot6\cdot6=\fbox{(E) 216}$.

 


	\item (\textit{AHSME 1986}) Let $p(x) = x^{2} + bx + c$, where $b$ and $c$ are integers. 
If $p(x)$ is a factor of both $x^{4} + 6x^{2} + 25$ and $3x^{4} + 4x^{2} + 28x + 5$, what is $p(1)$?


 $\textbf{(A)}\ 0\qquad \textbf{(B)}\ 1\qquad \textbf{(C)}\ 2\qquad \textbf{(D)}\ 4\qquad \textbf{(E)}\ 8$


 
 \textbf{Answer: }D



\textbf{Solution}

$p(x)$ must be a factor of $3(x^4+6x^2+25)-(3x^4+4x^2+28x+5)=14x^2-28x+70=14(x^2-2x+5)$.

 Therefore $p(x)=x^2 -2x+5$ and $p(1)=4$.

 The answer is $\fbox{(D) 4}$.

 


	\item (\textit{AHSME 1986}) If $\lfloor x\rfloor$ is the greatest integer less than or equal to $x$, then
$\sum_{N=1}^{1024} \lfloor \log_{2}N\rfloor =$


 $\textbf{(A)}\ 8192\qquad \textbf{(B)}\ 8204\qquad \textbf{(C)}\ 9218\qquad \textbf{(D)}\ \lfloor\log_{2}(1024!)\rfloor\qquad \textbf{(E)}\ \text{none of these}$


 
 \textbf{Answer: }B



\textbf{Solution}

Because $1 \le N \le 1024$, we have $0 \le \lfloor \log_{2}N\rfloor \le 10$. We count how many times $\lfloor \log_{2}N\rfloor$ attains a certain value. 

 For all $k$ except for $k=10$, we have that $\lfloor \log_{2}N\rfloor = k$ is satisfied by all $2^k \le N<2^{k+1}$, for a total of $2^k$ values of $N$. If $k=10$, $N$ can only have one value ($N=1024$). Thus, the desired sum looks like 
 \[\sum_{N=1}^{1024} \lfloor \log_{2}N\rfloor =1(0)+2(1)+4(2)+\dots+2^k(k)+\dots+2^{9}(9)+10\]

 Let $S$ be the desired sum without the $10$. 
 \[S=2(1)+4(2)+\dots+2^{9}(9)\] Multiplying by $2$ gives 
 \[2S=4(1)+8(2)+\dots+2^{10}(9)\] Subtracting the two equations gives 
 \[S=2^{10}(9)-(2+4+8+\dots+2^9)\] Summing the geometric sequence in parentheses and simplifying, we get 
 \[S=2^{10}(9)-2^{10}+2=2^{10}(8)+2=8194\] Finally, adding back the $10$ gives the desired answer $\fbox{(B) 8204}$

 


	\item (\textit{AHSME 1986}) It is desired to construct a right triangle in the coordinate plane so that its legs are parallel 
to the $x$ and $y$ axes and so that the medians to the midpoints of the legs lie on the lines $y = 3x + 1$
and $y = mx + 2$. The number of different constants $m$ for which such a triangle exists is


 $\textbf{(A)}\ 0\qquad \textbf{(B)}\ 1\qquad \textbf{(C)}\ 2\qquad \textbf{(D)}\ 3\qquad \textbf{(E)}\ \text{more than 3}$


 
 \textbf{Answer: }C



\textbf{Solution}

In \textit{any} right triangle with legs parallel to the axes, one median to the midpoint of a leg has slope $4$ times that of the other. This can easily be shown with coordinates: any triangle of this sort may be labelled with right angle at $P(a,b)$, other vertices $Q(a,b+2c)$ and $R(a-2d,b)$, and thus midpoints $(a,b+c)$ and $(a-d,b)$, so that the slopes are $\frac{c}{2d}$ and $\frac{2c}{d} = 4(\frac{c}{2d})$, thus showing that one is $4$ times the other as required.

 Thus in our problem, $m$ is either $3 \times 4 = 12$ or $3 \div 4 = \frac{3}{4}$. In fact, both are possible, and each for infinitely many triangles. We shall show this for $m=12$, and the argument is analogous for $m=\frac{3}{4}$. Take any right triangle with legs parallel to the axes and a hypotenuse with slope $12 \div 2 = 6$, e.g. the triangle with vertices $(0,0)$, $(1,0)$, and $(1,6)$. Then quick calculations show that the medians to the legs have slopes $12$ and $3$. Now translate the triangle (without rotating it) so that its medians intersect at the point where the lines $y=12x+2$ and $y=3x+1$ intersect. This forces the medians to lie on these lines (since their slopes are determined, and now we force them to go through a particular point; a slope and a point uniquely determine a line). Finally, for any central dilation of this triangle (a larger or smaller triangle with the same centroid and sides parallel to this one's sides), the medians will still lie on these lines, showing the "infinitely many" part of the result.

 Hence, to sum up, $m$ can in fact be both $12$ or $\frac{3}{4}$, which is exactly $2$ values, i.e. $\boxed{C}$.

 


	\item (\textit{AHSME 1986}) In the adjoining figure, $AB$ is a diameter of the circle, $CD$ is a chord parallel to $AB$, 
and $AC$ intersects $BD$ at $E$, with $\angle AED = \alpha$. The ratio of the area of $\triangle CDE$ to that of $\triangle ABE$ is  


 
                    \begin{figure}[H]
                        centering
                        \includegraphics[width=0.5\textwidth]{Problem/AHSME_I_1986_27_Problem_Figure_1.png}
                    \end{figure}
                    


 $\textbf{(A)}\ \cos\ \alpha\qquad \textbf{(B)}\ \sin\ \alpha\qquad \textbf{(C)}\ \cos^2\alpha\qquad \textbf{(D)}\ \sin^2\alpha\qquad \textbf{(E)}\ 1-\sin\ \alpha$


 
 \textbf{Answer: }C



\textbf{Solution}

$ABE$ and $DCE$ are similar isosceles triangles. It remains to find the square of the ratio of their sides. Draw in $AD$. Because $AB$ is a diameter, $\angle ADB=\angle ADE=90^{\circ}$. Thus, 
 \[\frac{DE}{AE}=\cos\alpha\] So 
 \[\frac{DE^2}{AE^2}=\cos^2\alpha\] The answer is thus $\fbox{(C)}$.

 


	\item (\textit{AHSME 1986}) $ABCDE$ is a regular pentagon. $AP, AQ$ and $AR$ are the perpendiculars dropped from $A$ onto $CD, CB$ extended and $DE$ extended, 
respectively. Let $O$ be the center of the pentagon. If $OP = 1$, then $AO + AQ + AR$ equals


 
                    \begin{figure}[H]
                        centering
                        \includegraphics[width=0.5\textwidth]{Problem/AHSME_I_1986_28_Problem_Figure_1.png}
                    \end{figure}
                    


 $\textbf{(A)}\ 3\qquad \textbf{(B)}\ 1 + \sqrt{5}\qquad \textbf{(C)}\ 4\qquad \textbf{(D)}\ 2 + \sqrt{5}\qquad \textbf{(E)}\ 5$


 
 \textbf{Answer: }C



\textbf{Solution 1}

To solve the problem, we compute the area of regular pentagon $ABCDE$ in two different ways. First, we can divide regular pentagon $ABCDE$ into five congruent triangles.
                    \begin{figure}[H]
                        centering
                        \includegraphics[width=0.5\textwidth]{Solution/AHSME_I_1986_28_Solution_Figure_1.png}
                    \end{figure}
                    

 If $s$ is the side length of the regular pentagon, then each of the triangles $AOB$, $BOC$, $COD$, $DOE$, and $EOA$ has base $s$ and height 1, so the area of regular pentagon $ABCDE$ is $5s/2$.

 Next, we divide regular pentagon $ABCDE$ into triangles $ABC$, $ACD$, and $ADE$.

 
                    \begin{figure}[H]
                        centering
                        \includegraphics[width=0.5\textwidth]{Solution/AHSME_I_1986_28_Solution_Figure_2.png}
                    \end{figure}
                    Triangle $ACD$ has base $s$ and height $AP = AO + 1$. Triangle $ABC$ has base $s$ and height $AQ$. Triangle $ADE$ has base $s$ and height $AR$. Therefore, the area of regular pentagon $ABCDE$ is also
 \[\frac{s}{2} (AO + AQ + AR + 1).\]Hence,
 \[\frac{s}{2} (AO + AQ + AR + 1) = \frac{5s}{2},\]which means $AO + AQ + AR + 1 = 5$, or $AO + AQ + AR = \boxed{4}$. The answer is $\boxed{(C)}$.

 



\textbf{Solution 2}

Now, we know that angle $D$ has measure $\frac{180 \cdot 3}{5} = 108$. Since
 \[\sin 54 = \frac{OP}{DO} = \frac{1}{DO}, DO = \frac{1}{\sin 54}\]
 \[\tan 54 = \frac{OP}{DP} = \frac{1}{DP}, DP = \frac{1}{\tan 54}\]Therefore, $AB = 2DP = \frac{2}{\tan 54}$.
 \[\sin 72 = \frac{AQ}{AB} = AQ \tan 54 \cdot \frac{1}{2}, AQ = \frac{2 \sin 72}{\tan 54}\]Therefore, $AO + AQ + AR = AO + 2AQ = \frac{1}{\sin 54}+\frac{4 \sin 72}{\tan 54} = \frac{1}{\sin 54} + 8 \sin 36 \cos 54 = \frac{1}{\cos 36} + 8-8\cos^2(36)$. Recalling that $\cos 36 = \frac{1 + \sqrt{5}}{4}$ gives a final answer of $\boxed{4}$.

 


	\item (\textit{AHSME 1986}) Two of the altitudes of the scalene triangle $ABC$ have length $4$ and $12$. 
If the length of the third altitude is also an integer, what is the biggest it can be?


 $\textbf{(A)}\ 4\qquad \textbf{(B)}\ 5\qquad \textbf{(C)}\ 6\qquad \textbf{(D)}\ 7\qquad \textbf{(E)}\ \text{none of these}$


 
 \textbf{Answer: }B



\textbf{Solution 1}

Assume we have a scalene triangle $ABC$. Arbitrarily, let $12$ be the height to base $AB$ and $4$ be the height to base $AC$. Due to area equivalences, the base $AC$ must be three times the length of $AB$. 

 Let the base $AB$ be $x$, thus making $AC = 3x$. Thus, setting the final height to base $BC$ to $h$, we note that (by area equivalence) $\frac{BC \cdot h}{2} = \frac{3x \cdot 4}{2} = 6x$. Thus, $h = \frac{12x}{BC}$. We note that to maximize $h$ we must minimize $BC$. Using the triangle inequality, $BC + AB > AC$, thus $BC + x > 3x$ or $BC > 2x$. The minimum value of $BC$ is $2x$, which would output $h = 6$. However, because $BC$ must be larger than $2x$, the minimum integer height must be $5$.

 $\fbox{B}$

 


 



\textbf{Solution 2}

The reciprocals of the altitudes of a triangle themselves form a triangle - this can be easily proven. Let our desired altitude be $a$. 

 We have $\frac{1}{a}<\frac{1}{4}+\frac{1}{12}=\frac{1}{3}$, which implies $a>3$. We also have $\frac{1}{a}>\frac{1}{4}-\frac{1}{12}=\frac{1}{6}$, which implies $a<6$. Therefore the maximum integral value of $a$ is 5.

 $\fbox{B}$.

 


	\item (\textit{AHSME 1986}) The number of real solutions $(x,y,z,w)$ of the simultaneous equations
$2y = x + \frac{17}{x}, 2z = y + \frac{17}{y}, 2w = z + \frac{17}{z}, 2x = w + \frac{17}{w}$
is


 $\textbf{(A)}\ 1\qquad \textbf{(B)}\ 2\qquad \textbf{(C)}\ 4\qquad \textbf{(D)}\ 8\qquad \textbf{(E)}\ 16$


 
 \textbf{Answer: }B



\textbf{Solution 1}

Consider the cases $x>0$ and $x<0$, and also note that by AM-GM, for any positive number $a$, we have $a+\frac{17}{a} \geq 2\sqrt{17}$, with equality only if $a = \sqrt{17}$. Thus, if $x>0$, considering each equation in turn, we get that $y \geq \sqrt{17}, z \geq \sqrt{17}, w \geq \sqrt{17}$, and finally $x \geq \sqrt{17}$. 

 Now suppose $x > \sqrt{17}$. Then $y - \sqrt{17} = \frac{x^{2}+17}{2x} - \sqrt{17} = (\frac{x-\sqrt{17}}{2x})(x-\sqrt{17}) < \frac{1}{2}(x-\sqrt{17})$, so that $x > y$. Similarly, we can get $y > z$, $z > w$, and $w > x$, and combining these gives $x > x$, an obvious contradiction.

 Thus we must have $x \geq \sqrt{17}$, but $x \ngtr \sqrt{17}$, so if $x > 0$, the only possibility is $x = \sqrt{17}$, and analogously from the other equations we get $x = y = z = w = \sqrt{17}$; indeed, by substituting, we verify that this works.

 As for the other case, $x < 0$, notice that $(x,y,z,w)$ is a solution if and only if $(-x,-y,-z,-w)$ is a solution, since this just negates both sides of each equation and so they are equivalent. Thus the only other solution is $x = y = z = w = -\sqrt{17}$, so that we have $2$ solutions in total, and therefore the answer is $\boxed{B}$.

 



\textbf{Solution 2}

Let $(a, b, c, d) = \left(\frac{x}{\sqrt{17}}, \frac{y}{\sqrt{17}}, \frac{z}{\sqrt{17}}, \frac{w}{\sqrt{17}}\right)$.  Then these equations become $2b = a + \frac{1}{a}$, $2c = b + \frac{1}{b}$, $2d = c + \frac{1}{c}$, and $a = 2d + \frac{1}{d}$ after substituting and cancelling.  Notice that $(a, b, c, d)$ must either all be positive or all be negative, and that $(a, b, c, d)$ is a solution if and only if $(-a, -b, -c, -d)$ is a solution.

 Suppose $(a, b, c, d)$ are all positive.  Then by AM-GM, $2a = d + \frac{1}{d} \geq 2\sqrt{d \cdot \frac{1}{d}} = 2$, so $a \geq 1$.  By similar logic, $b \geq 1$, $c \geq 1$, and $d \geq 1$.  Therefore, $a \geq \frac{1}{a}$, $b \geq \frac{1}{b}$, $c \geq \frac{1}{c}$, and $d \geq \frac{1}{d}$.

 So $2a = d + \frac{1}{d} \leq 2d = c + \frac{1}{c} \leq 2c = b + \frac{1}{b} \leq 2b = a + \frac{1}{a} \leq 2a$.  Thus $a \leq d \leq c \leq b \leq a$, so $a = b = c = d$.  So $2a = a + \frac{1}{a}$, that is, $a = 1$ since $a$ is positive.  Therefore, the only positive solution is $(a, b, c, d) = (1, 1, 1, 1)$, and thus the only negative solution is $(a, b, c, d) = (-1, -1, -1, -1)$.

 Since the only solutions for $(a, b, c, d)$ are $(1, 1, 1, 1)$ and $(-1, -1, -1, -1)$, the only $\boxed{(\mathbf{B})\ 2}$ solutions for $(x, y, z, w)$ are $(\sqrt{17}, \sqrt{17}, \sqrt{17}, \sqrt{17})$ and $(-\sqrt{17}, -\sqrt{17}, -\sqrt{17}, -\sqrt{17})$.

 -j314andrews

 


\end{enumerate}